% !TeX root = ../thuthesis-example.tex

\chapter{功能设计}
\section{需求分析}
\subsection{功能性需求分析}
功能性需求是根据平台的核心业务流程和用户交互场景得出的，描述了系统必须为用户提供的具体服务和功能。首先，系统必须提供完整的用户账户管理功能，支持学习者注册、登录并创建包含目标方言选择的个人档案。在内容呈现上，平台需要以可视化的知识节点网络作为主学习界面，引导用户进行结构化学习，同时还应提供一个可自由探索的交互式知识图谱，以满足用户对知识进行溯源和关联发现的需求。当用户点击任一学习单元时，系统必须能展示其详细内容，并提供标准发音的音频播放功能。

平台的核心功能需求围绕着一个“学-练-评-荐”的智能闭环展开。系统需支持用户录制发音，并能对其进行实时的声学分析。关键需求在于，系统必须基于知识图谱进行多维度的智能诊断，不仅判断对错，还要提供具体、可解释的改进建议。练习完成后，系统必须能自动生成专属的“错题本”，并基于对错题记录和用户历史表现的综合分析，动态地为用户规划和推荐下一步的学习路径，从而实现真正的个性化指导。

\subsection{非功能性需求分析}
非功能性需求是衡量平台质量与用户体验的关键，涉及性能、可靠性、易用性等多个方面。在性能与准确性方面，平台必须满足严苛的要求。发音诊断与反馈的响应时间需控制在数秒以内，以保障交互的流畅性，同时后台需具备支持多人并发学习的能力。更重要的是，作为平台的核心价值，发音诊断的准确率必须达到高可信水平，这是赢得学习者信任的基石。此外，系统整体应保持高可用性，确保服务的稳定可靠，并保证用户数据在不同终端上的一致性。

在可用性层面，平台界面设计需力求简洁直观，确保不同年龄段的用户都能轻松上手；智能体给出的诊断建议也必须通俗易懂，避免使用过多的专业术语，使其具有真正的指导价值。从长远来看，系统的可扩展性与可维护性至关重要。平台架构必须支持未来便捷地增添新的方言种类，其模块化的设计也应确保各个功能组件可以独立升级和维护，以适应技术的持续发展。

\section{平台功能设计}

本项目的最终目标是落地为一个功能完备、体验流畅、且具备高度智能化的方言学习平台。平台的功能设计遵循以学习者为中心、以交互式练习为核心、以自适应规划为导向的原则，旨在将底层的复杂算法无缝转化为直观、高效的学习体验。整个平台的功能体系主要由以下几个核心模块构成。

\subsection{用户管理与学习目标设定模块}
这是用户与平台交互的入口。首次使用时，用户需进行注册并创建个人学习档案。本模块最核心的功能是多方言选择系统。用户可以从平台支持的方言库中（如粤语、上海话、闽南语等）选择自己希望学习的目标方言。一旦选定，平台的所有后续内容，包括知识呈现、发音练习、路径规划等，都将围绕这一目标方言展开。该模块负责持续记录用户的基本信息、学习方言种类以及整体学习进度。

\subsection{可视化知识节点学习模块}
当用户选定方言进入主学习界面后，平台将呈现一个基于知识图谱的可视化学习网络。在这个网络中，方言的知识体系被拆解成一个个相互关联的“学习节点”。每个节点代表一个具体的学习单元，可能是一个核心声母/韵母、一个独特的声调模式、一组易混淆的词汇，或一个特定的语法现象。用户可以根据引导或自主选择点击不同的节点，进入针对性的练习环节。这种设计将抽象的语言知识具象化，使用户能够清晰地看到自己的学习版图和知识掌握情况。

\subsection{交互式知识图谱探索模块}
为满足学习者对知识进行溯源和关联性探索的需求，并增强平台教学的透明度，本项目将设计一个交互式的知识图谱探索模块。该模块将底层的“方言声学知识图谱”以图形化、可交互的网络形式完整地呈现给用户。学习者不再局限于预设的学习节点，而是可以自由地在图谱中进行漫游、缩放和点击。例如，用户可以点击一个词汇节点，查看其连接的所有相关信息，包括其标准发音音频、所属的音系规则、以及具有相同发音特征的其他词汇等。此外，在“核心发音练习模块”中，当智能体给出诊断建议时，会提供直接跳转到知识图谱中对应节点的链接，方便学习者即时查看该知识点的完整上下文，从而实现从“被动纠错”到“主动探索”的升华。

\subsection{核心发音练习与智能诊断模块}
这是平台最核心的互动学习环节，也是“声学多模态认知校准框架”和“音系知识自进化算法”的主要应用场景。当用户进入一个学习节点后，智能体首先会从“方言声学知识图谱”中呈现相应的练习任务，并提供与图谱关联的真人标准发音作为示范。学习者在收听后，可以录制自己的发音，系统则会实时捕捉语音信号并进行声学分析，例如生成用于对比的语图。随后，作为本模块的核心，智能诊断引擎会将用户的发音特征与知识图谱中的标准模式进行精确比对，并给出多维度的、可解释的诊断报告。这种报告不仅判断对错，更能明确指出问题所在（如声调、声母或韵母），并提供具体的改进建议，从而构成一个完整的“学-练-评”闭环。

\subsection{自适应评估与学习路径规划模块}
该模块是实现个性化教学的关键，也是“个性化自适应学习算法”的具体体现。当用户完成一个节点的练习后，系统会自动将其发音不佳或回答错误的条目，连同智能体的诊断建议，汇集成一本专属的、可视化的“错题本”。这本错题本不仅是供用户随时回顾复习的工具，更是自适应路径规划的核心依据。智能体将综合分析当前节点的错题内容及用户的历史学习数据，动态更新该用户的“学习者状态模型”。基于此模型，系统最终会为用户推荐或规划出下一阶段最需巩固的知识节点，从而实现一条真正动态、智能的学习路径。